% ============================================================================
% DM-2026-006: NASA CUI Directive Requirements Absorption
% ============================================================================
% Thin wrapper for Decision Memorandum - see _template.tex for shared preamble
% ============================================================================

% ----------------------------------------------------------------------------
% DOCUMENT VARIABLES
% ----------------------------------------------------------------------------
\newcommand{\UniqueID}{DM-2026-006}
\newcommand{\DocumentDate}{February 17, 2026}
\newcommand{\AuthorName}{Bruce Dombrowski}
\newcommand{\AuthorTitle}{Systems Engineer}
\newcommand{\ToField}{V\&V Framework Stakeholders}
\newcommand{\SubjectField}{NASA CUI Directive Requirements Absorption}
\newcommand{\OPRField}{Compliance-Marking}

% ----------------------------------------------------------------------------
% DOCUMENT CONTENT
% ----------------------------------------------------------------------------
\newcommand{\dmContent}{%
\begin{enumerate}

\dmsection{Purpose}

This memorandum documents the decision to absorb NASA's CUI directive hierarchy into formal requirements artifacts. Two documents are captured: the current active NPR 2810.7 ``Controlled Unclassified Information'' and the expired interim NID 2810.135. Together they demonstrate the systems-engineering framework's ability to absorb external governing-authority requirements and produce formal Verification \& Validation (V\&V) artifacts.

\dmsection{Background}

NASA's Controlled Unclassified Information (CUI) program standardizes how sensitive information is marked, handled, and shared across the agency. The program is governed by a directive hierarchy:

\begin{itemize}
    \item \textbf{Federal Level:} Executive Order 13556 (establishing the CUI program) and 32 CFR Part 2002 (implementation regulation)
    \item \textbf{NASA Policy:} NPD 2810.1F ``NASA Information Security Policy'' (effective Jan 21, 2022; expires Jan 21, 2027) --- establishes overarching information security policy and mandates the CUI program per EO 13556
    \item \textbf{NASA Procedure (Current):} NPR 2810.7 ``Controlled Unclassified Information'' Change 1 (effective Oct 22, 2021; expires Oct 22, 2026) --- 23 subsections covering all CUI management procedures. This is the \textbf{current active governing document} for production use.
    \item \textbf{NASA Interim (Expired):} NID 2810.135 ``Controlled Unclassified Information'' (effective Feb 2, 2021; expired Feb 2, 2022) --- interim directive whose requirements were carried forward into NPR 2810.7
\end{itemize}

NID 2810.135 is publicly available at: \url{https://nodis3.gsfc.nasa.gov/OPD_docs/NID_2810_135_.pdf}

NPR 2810.7 is available at: \url{https://nodis3.gsfc.nasa.gov/displayDir.cfm?t=NPR&c=2810&s=7}

\dmsection{Analysis}

\textbf{Why use expired NID 2810.135 rather than active NPR 2810.7?}

\begin{itemize}
    \item \textbf{Safe test data:} The expired status means no compliance obligations arise from citing it. It cannot be mistaken for an active requirement source.
    \item \textbf{Framework demonstration:} The goal is to prove the REQ/VER toolchain can absorb a complex, multi-chapter directive with diverse requirement types (technical, organizational, informational). NID 2810.135 serves this purpose.
    \item \textbf{Full lifecycle visibility:} Using an expired directive shows the complete lifecycle: active $\rightarrow$ superseded $\rightarrow$ archived. The successor relationship (NID 2810.135 $\rightarrow$ NPR 2810.7) is documented in the artifact metadata.
    \item \textbf{Requirements equivalence:} All normative requirements in NID 2810.135 were carried forward into NPR 2810.7. The test data is representative of real CUI requirements without creating ambiguity about which standard is authoritative.
    \item \textbf{Publicly shareable:} NID 2810.135 is a public document. Artifacts derived from it can be freely shared, discussed, and used as examples.
\end{itemize}

\dmsection{Decision}

\textbf{Approved:} Absorb both NASA CUI directives as formal requirements artifacts.

\textbf{Artifacts to be produced:}
\begin{itemize}
    \item \textbf{REQ-2026-002:} JSON requirements extraction from NID 2810.135 (expired interim directive) --- serves as framework test data demonstrating full lifecycle
    \item \textbf{REQ-2026-003:} JSON requirements extraction from NPR 2810.7 (current active procedural requirement) --- the authoritative CUI standard for production compliance
    \item \textbf{VER-2026-002:} LaTeX verification document verifying ORBIT-RIFD-Status against applicable technical requirements from both REQ documents
    \item \textbf{build-requirements-pdf.py:} Parameterized build script generating PDF from JSON (adapted from SendCUIEmail REQ-2026-001 toolchain)
\end{itemize}

\textbf{Scope:} REQ extraction from both directives, plus VER against ORBIT-RIFD-Status as a real-world verification target. ORBIT-RIFD-Status is ideal because it handles CUI data daily (ISS program CEFs/PIAs), implements CUI markings on outputs, uses DPAPI database encryption, and requires CAC/PIV authentication.

\textbf{Rationale for both documents:}
\begin{itemize}
    \item \textbf{NPR 2810.7} is the current governing authority --- production projects reference REQ-2026-003
    \item \textbf{NID 2810.135} (expired) demonstrates the expired-directive lifecycle and serves as safe, shareable test data
    \item Requirements are substantively equivalent; the interim directive's requirements were carried forward into NPR 2810.7
\end{itemize}

\dmsection{Implementation Notes}

\begin{itemize}
    \item Both REQ JSONs include a \texttt{source\_document} metadata block recording: document ID, status, successor/predecessor references, and provenance URLs
    \item REQ-2026-002 (NID) includes an expired-status warning and test-data disclaimer
    \item REQ-2026-003 (NPR) is marked as the current active governing document
    \item Each requirement carries an \texttt{applicability} field: \texttt{technical} (verifiable against software), \texttt{organizational} (requires process/policy), or \texttt{informational} (reference material)
    \item The build script is backward-compatible with REQ-2026-001 (graceful fallback when new fields are absent)
    \item VER-2026-002 lives in the ORBIT-CEF-Status repository (co-located with the implementation it verifies)
\end{itemize}

\dmsection{References}

\begin{itemize}
    \item NID 2810.135: \url{https://nodis3.gsfc.nasa.gov/OPD_docs/NID_2810_135_.pdf}
    \item NPR 2810.7: \url{https://nodis3.gsfc.nasa.gov/displayDir.cfm?t=NPR&c=2810&s=7}
    \item NPD 2810.1F: \url{https://nodis3.gsfc.nasa.gov/displayDir.cfm?t=NPD&c=2810&s=1E}
    \item NASA CUI Program: \url{https://www.nasa.gov/controlled-unclassified-information-cui-and-why-it-matters/}
    \item EO 13556: \url{https://www.archives.gov/cui/registry/policy-guidance}
    \item 32 CFR Part 2002: \url{https://www.ecfr.gov/current/title-32/subtitle-B/chapter-XX/part-2002}
    \item NARA CUI Registry: \url{https://www.archives.gov/cui}
    \item REQ-2026-001 (SendCUIEmail precedent): Cryptographic Compliance Requirements
    \item VER-2026-001 (SendCUIEmail precedent): Cryptographic Compliance Verification
    \item Related Decision: DM-2026-001 (SF901 LaTeX Recreation)
\end{itemize}

\end{enumerate}
}

% ----------------------------------------------------------------------------
% INCLUDE SHARED TEMPLATE
% ----------------------------------------------------------------------------
% ============================================================================
% Decision Memorandum Base Template
% ============================================================================
% This is the shared base template for all Decision Memoranda.
% DO NOT edit this file for individual decisions.
%
% Usage: Create a thin wrapper that defines variables and content, then:
%   % ============================================================================
% Decision Memorandum Base Template
% ============================================================================
% This is the shared base template for all Decision Memoranda.
% DO NOT edit this file for individual decisions.
%
% Usage: Create a thin wrapper that defines variables and content, then:
%   % ============================================================================
% Decision Memorandum Base Template
% ============================================================================
% This is the shared base template for all Decision Memoranda.
% DO NOT edit this file for individual decisions.
%
% Usage: Create a thin wrapper that defines variables and content, then:
%   \input{_template.tex}
%
% Required variables (define before \input):
%   \UniqueID       - e.g., DM-2026-001
%   \DocumentDate   - e.g., January 19, 2026
%   \AuthorName     - Signer name
%   \AuthorTitle    - Signer title
%   \ToField        - Distribution
%   \SubjectField   - Subject line
%   \OPRField       - Office of Primary Responsibility
%   \dmContent      - The full content (enumerate environment with \dmsection items)
%
% ============================================================================

\documentclass[11pt,letterpaper]{article}

% ============================================================================
% PACKAGES
% ============================================================================
\usepackage[margin=1in, top=1.15in, headheight=30pt]{geometry}
\usepackage{graphicx}
\usepackage{fancyhdr}
\usepackage{lastpage}
\usepackage{enumitem}
\usepackage{xcolor}
\usepackage{hyperref}
\usepackage{tabularx}

% ============================================================================
% DOCUMENT CONFIGURATION
% ============================================================================
\hypersetup{
    colorlinks=true,
    linkcolor=blue,
    urlcolor=blue,
    pdfborder={0 0 0}
}

% Remove paragraph indentation
\setlength{\parindent}{0pt}
\setlength{\parskip}{6pt}

% List formatting - match Word style
\setlist[enumerate]{leftmargin=0.5in, labelwidth=0.25in, labelsep=0.1in, align=left, itemsep=12pt, topsep=12pt}
\setlist[enumerate,1]{label=\arabic*.}
\setlist[itemize]{leftmargin=0.5in, itemsep=3pt, topsep=6pt}

% ============================================================================
% HEADER AND FOOTER CONFIGURATION
% ============================================================================
\pagestyle{fancy}
\fancyhf{}

% Header: text-only, Boeing / ISS Program identification
\fancyhead[L]{\large\textbf{Decision Memorandum}\\[-2pt]\scriptsize Boeing \textbar{} International Space Station Program}
\fancyhead[R]{\small \UniqueID\\[-2pt]\scriptsize \DocumentDate}
\renewcommand{\headrulewidth}{0.4pt}

% Footer: Unique ID (left), Page X of Y (center), Date (right)
\fancyfoot[L]{\small \UniqueID}
\fancyfoot[C]{\small Page \thepage\ of \pageref{LastPage}}
\fancyfoot[R]{\small \DocumentDate}
\renewcommand{\footrulewidth}{0.4pt}

% ============================================================================
% CUSTOM COMMANDS
% ============================================================================
% Bold section headers for numbered sections
\newcommand{\dmsection}[1]{\item \textbf{#1}\par}

% ============================================================================
% BEGIN DOCUMENT
% ============================================================================
\begin{document}

% ----------------------------------------------------------------------------
% UNIQUE ID (top right)
% ----------------------------------------------------------------------------
\hfill \UniqueID

\vspace{0.25in}

% ----------------------------------------------------------------------------
% SIGNATURE BLOCK
% ----------------------------------------------------------------------------
\underline{\hspace{2.5in}}\\[2pt]
\AuthorName\\
\AuthorTitle

\vspace{0.25in}

% ----------------------------------------------------------------------------
% MEMO HEADER FIELDS
% ----------------------------------------------------------------------------
\textbf{DATE:} \DocumentDate

\textbf{TO:} \ToField

\textbf{SUBJECT:} \SubjectField

\textbf{OFFICE OF PRIMARY RESPONSIBILITY:} \OPRField

\vspace{0.25in}

% ============================================================================
% MAIN CONTENT - NUMBERED SECTIONS
% ============================================================================
\dmContent

\end{document}

%
% Required variables (define before \input):
%   \UniqueID       - e.g., DM-2026-001
%   \DocumentDate   - e.g., January 19, 2026
%   \AuthorName     - Signer name
%   \AuthorTitle    - Signer title
%   \ToField        - Distribution
%   \SubjectField   - Subject line
%   \OPRField       - Office of Primary Responsibility
%   \dmContent      - The full content (enumerate environment with \dmsection items)
%
% ============================================================================

\documentclass[11pt,letterpaper]{article}

% ============================================================================
% PACKAGES
% ============================================================================
\usepackage[margin=1in, top=1.15in, headheight=30pt]{geometry}
\usepackage{graphicx}
\usepackage{fancyhdr}
\usepackage{lastpage}
\usepackage{enumitem}
\usepackage{xcolor}
\usepackage{hyperref}
\usepackage{tabularx}

% ============================================================================
% DOCUMENT CONFIGURATION
% ============================================================================
\hypersetup{
    colorlinks=true,
    linkcolor=blue,
    urlcolor=blue,
    pdfborder={0 0 0}
}

% Remove paragraph indentation
\setlength{\parindent}{0pt}
\setlength{\parskip}{6pt}

% List formatting - match Word style
\setlist[enumerate]{leftmargin=0.5in, labelwidth=0.25in, labelsep=0.1in, align=left, itemsep=12pt, topsep=12pt}
\setlist[enumerate,1]{label=\arabic*.}
\setlist[itemize]{leftmargin=0.5in, itemsep=3pt, topsep=6pt}

% ============================================================================
% HEADER AND FOOTER CONFIGURATION
% ============================================================================
\pagestyle{fancy}
\fancyhf{}

% Header: text-only, Boeing / ISS Program identification
\fancyhead[L]{\large\textbf{Decision Memorandum}\\[-2pt]\scriptsize Boeing \textbar{} International Space Station Program}
\fancyhead[R]{\small \UniqueID\\[-2pt]\scriptsize \DocumentDate}
\renewcommand{\headrulewidth}{0.4pt}

% Footer: Unique ID (left), Page X of Y (center), Date (right)
\fancyfoot[L]{\small \UniqueID}
\fancyfoot[C]{\small Page \thepage\ of \pageref{LastPage}}
\fancyfoot[R]{\small \DocumentDate}
\renewcommand{\footrulewidth}{0.4pt}

% ============================================================================
% CUSTOM COMMANDS
% ============================================================================
% Bold section headers for numbered sections
\newcommand{\dmsection}[1]{\item \textbf{#1}\par}

% ============================================================================
% BEGIN DOCUMENT
% ============================================================================
\begin{document}

% ----------------------------------------------------------------------------
% UNIQUE ID (top right)
% ----------------------------------------------------------------------------
\hfill \UniqueID

\vspace{0.25in}

% ----------------------------------------------------------------------------
% SIGNATURE BLOCK
% ----------------------------------------------------------------------------
\underline{\hspace{2.5in}}\\[2pt]
\AuthorName\\
\AuthorTitle

\vspace{0.25in}

% ----------------------------------------------------------------------------
% MEMO HEADER FIELDS
% ----------------------------------------------------------------------------
\textbf{DATE:} \DocumentDate

\textbf{TO:} \ToField

\textbf{SUBJECT:} \SubjectField

\textbf{OFFICE OF PRIMARY RESPONSIBILITY:} \OPRField

\vspace{0.25in}

% ============================================================================
% MAIN CONTENT - NUMBERED SECTIONS
% ============================================================================
\dmContent

\end{document}

%
% Required variables (define before \input):
%   \UniqueID       - e.g., DM-2026-001
%   \DocumentDate   - e.g., January 19, 2026
%   \AuthorName     - Signer name
%   \AuthorTitle    - Signer title
%   \ToField        - Distribution
%   \SubjectField   - Subject line
%   \OPRField       - Office of Primary Responsibility
%   \dmContent      - The full content (enumerate environment with \dmsection items)
%
% ============================================================================

\documentclass[11pt,letterpaper]{article}

% ============================================================================
% PACKAGES
% ============================================================================
\usepackage[margin=1in, top=1.15in, headheight=30pt]{geometry}
\usepackage{graphicx}
\usepackage{fancyhdr}
\usepackage{lastpage}
\usepackage{enumitem}
\usepackage{xcolor}
\usepackage{hyperref}
\usepackage{tabularx}

% ============================================================================
% DOCUMENT CONFIGURATION
% ============================================================================
\hypersetup{
    colorlinks=true,
    linkcolor=blue,
    urlcolor=blue,
    pdfborder={0 0 0}
}

% Remove paragraph indentation
\setlength{\parindent}{0pt}
\setlength{\parskip}{6pt}

% List formatting - match Word style
\setlist[enumerate]{leftmargin=0.5in, labelwidth=0.25in, labelsep=0.1in, align=left, itemsep=12pt, topsep=12pt}
\setlist[enumerate,1]{label=\arabic*.}
\setlist[itemize]{leftmargin=0.5in, itemsep=3pt, topsep=6pt}

% ============================================================================
% HEADER AND FOOTER CONFIGURATION
% ============================================================================
\pagestyle{fancy}
\fancyhf{}

% Header: text-only, Boeing / ISS Program identification
\fancyhead[L]{\large\textbf{Decision Memorandum}\\[-2pt]\scriptsize Boeing \textbar{} International Space Station Program}
\fancyhead[R]{\small \UniqueID\\[-2pt]\scriptsize \DocumentDate}
\renewcommand{\headrulewidth}{0.4pt}

% Footer: Unique ID (left), Page X of Y (center), Date (right)
\fancyfoot[L]{\small \UniqueID}
\fancyfoot[C]{\small Page \thepage\ of \pageref{LastPage}}
\fancyfoot[R]{\small \DocumentDate}
\renewcommand{\footrulewidth}{0.4pt}

% ============================================================================
% CUSTOM COMMANDS
% ============================================================================
% Bold section headers for numbered sections
\newcommand{\dmsection}[1]{\item \textbf{#1}\par}

% ============================================================================
% BEGIN DOCUMENT
% ============================================================================
\begin{document}

% ----------------------------------------------------------------------------
% UNIQUE ID (top right)
% ----------------------------------------------------------------------------
\hfill \UniqueID

\vspace{0.25in}

% ----------------------------------------------------------------------------
% SIGNATURE BLOCK
% ----------------------------------------------------------------------------
\underline{\hspace{2.5in}}\\[2pt]
\AuthorName\\
\AuthorTitle

\vspace{0.25in}

% ----------------------------------------------------------------------------
% MEMO HEADER FIELDS
% ----------------------------------------------------------------------------
\textbf{DATE:} \DocumentDate

\textbf{TO:} \ToField

\textbf{SUBJECT:} \SubjectField

\textbf{OFFICE OF PRIMARY RESPONSIBILITY:} \OPRField

\vspace{0.25in}

% ============================================================================
% MAIN CONTENT - NUMBERED SECTIONS
% ============================================================================
\dmContent

\end{document}

